% ============================================================================
% APPENDIX BE: LIT-ENKE - Moral Universalism & Cognitive Uncertainty
% ============================================================================
% Category: Literature Integration (LIT-)
% Language: English
% EBF Integration: Culture, cognition, and moral preferences
% ============================================================================

\chapter*{Appendix UNMAPPED_BE: Benjamin Enke Research Integration}
\addcontentsline{toc}{chapter}{Appendix UNMAPPED_BE: LIT-ENKE - Moral Universalism \& Cognitive Uncertainty}
\label{app:enke}

% ============================================================================
% HEADER BLOCK
% ============================================================================
\begin{tcolorbox}[colback=blue!5!white, colframe=blue!75!black, title=Appendix Metadata]
\textbf{Code:} BE \\
\textbf{Category:} LIT- (Literature Integration) \\
\textbf{Title:} Moral Universalism \& Cognitive Uncertainty \\
\textbf{Version:} 1.0 \\
\textbf{Status:} Complete \\
\textbf{Primary Author Focus:} Benjamin Enke \\
\textbf{Papers Covered:} 20 \\
\textbf{Dependencies:} CORE-AWARE (AU), CONTEXT-CULTURE, LIT-FALK (AE), CORE-WHEN (V) \\
\textbf{EBF Connections:} Cultural Ψ parameters, cognitive noise, moral preferences
\end{tcolorbox}

% ============================================================================
% CROSS-REFERENCE MAP
% ============================================================================
\begin{tcolorbox}[colback=green!5!white, colframe=green!75!black, title=Cross-Reference Map]
\textbf{Outgoing References:}
\begin{itemize}[nosep]
    \item CORE-AWARE (AU): Cognitive uncertainty $\rightarrow$ awareness limitations
    \item CORE-WHEN (V): Complexity-driven time preferences
    \item CONTEXT-CULTURE: Moral universalism parameters
    \item LIT-FALK (AE): Global Preferences Survey collaboration
    \item LIT-GNEEZY (GN): Truth-telling and deception research
    \item LIT-GABAIX (XG): Cognitive noise and bounded rationality
\end{itemize}

\textbf{Incoming References:}
\begin{itemize}[nosep]
    \item Chapter 9 (Context): Cultural variation in Ψ
    \item Chapter 11 (Awareness): Cognitive uncertainty modeling
    \item PREDICT-03: Culture-based predictions
\end{itemize}
\end{tcolorbox}

% ============================================================================
% ABSTRACT
% ============================================================================
\section*{Abstract}

Benjamin Enke has made fundamental contributions to understanding how \textbf{culture}, \textbf{cognition}, and \textbf{moral values} shape economic behavior. His work bridges anthropology, psychology, and economics through rigorous empirical methods. This appendix integrates Enke's 20 key papers into the EBF framework, with particular emphasis on:
\begin{enumerate}
    \item \textbf{Moral Universalism vs. Communalism}: How kinship structures shape cooperation
    \item \textbf{Cognitive Uncertainty}: Noise in decision-making as a unifying framework
    \item \textbf{Correlation Neglect}: Systematic errors in belief formation
    \item \textbf{Complexity and Time}: Why complexity generates apparent hyperbolic discounting
    \item \textbf{Cultural Transmission}: How ancestral environments shape modern preferences
\end{enumerate}

% ============================================================================
% QUICK REFERENCE
% ============================================================================
\begin{tcolorbox}[colback=yellow!10!white, colframe=orange!75!black, title=Quick Reference: Key Enke Parameters]
\begin{tabular}{lll}
\textbf{Parameter} & \textbf{Value} & \textbf{Source} \\
\hline
Moral Universalism Index & 0-1 scale & Enke (2019, 2022) \\
Cognitive uncertainty $\sigma$ & 0.15-0.35 & Enke \& Graeber (2023) \\
Correlation neglect & 30-50\% & Enke \& Zimmermann (2019) \\
Complexity-time elasticity & 0.4-0.6 & Enke, Graeber \& Oprea (2024) \\
Kinship tightness effect & $\beta = 0.25$ & Enke (2019) \\
\end{tabular}
\end{tcolorbox}

% ============================================================================
% FUNDAMENTAL QUESTION
% ============================================================================
% -----------------------------------------------------------------------------
% APPENDIX SCOPE BOX
% -----------------------------------------------------------------------------

\begin{tcolorbox}[
    colback=orange!5!white,
    colframe=orange!75!black,
    title={\textbf{Appendix Scope: Ziel / In-Scope / Out-of-Scope / Constraints}},
    fonttitle=\bfseries
]

\textbf{Ziel (Objective):}
\begin{itemize}[nosep]
    \item How do deep cultural structures and cognitive limitations jointly determine economic preferences and behavior?
\end{itemize}

\textbf{In-Scope:}
\begin{itemize}[nosep]
    \item Fundamental Question
    \item Cognitive Uncertainty Framework
    \item Moral Values in Economic Contexts
    \item EBF Integration
    \item Worked Example: Predicting Cooperation Patterns
\end{itemize}

\textbf{Out-of-Scope (delegiert an andere Appendices):}
\begin{itemize}[nosep]
    \item Glossary $\rightarrow$ Appendix G
\end{itemize}

\textbf{Constraints (Anwendungsgrenzen):}
\begin{itemize}[nosep]
    \item [TODO: Define when this appendix does NOT apply]
    \item [TODO: Define required prerequisites]
    \item [TODO: Define known limitations]
\end{itemize}

\textbf{Lieferobjekte (Deliverables):}
\begin{itemize}[nosep]
    \item 1 Table(s)
    \item 6 Equation(s)
    \item Worked Example(s)
\end{itemize}

\end{tcolorbox}


\section{Fundamental Question}

\begin{tcolorbox}[colback=red!5!white, colframe=red!75!black, title=Central Question]
\textbf{How do deep cultural structures and cognitive limitations jointly determine economic preferences and behavior?}
\end{tcolorbox}

Enke's answer: Cultural evolution (shaped by kinship systems and ancestral environments) creates persistent differences in moral values, while cognitive noise introduces systematic biases that interact with complexity to generate behavioral anomalies.

% ============================================================================
% THEORETICAL FOUNDATIONS
% ============================================================================
\section{Theoretical Foundations}

\subsection{Moral Universalism vs. Communalism}

Enke (2019) develops a framework distinguishing two moral orientations:

\begin{definition}[Moral Universalism]
A moral system where ethical obligations apply \textbf{universally} to all humans, regardless of group membership. Associated with:
\begin{itemize}
    \item Impartial institutions
    \item Impersonal exchange
    \item Democratic governance
\end{itemize}
\end{definition}

\begin{definition}[Moral Communalism]
A moral system where ethical obligations are \textbf{particularistic}, varying by relationship type (in-group vs. out-group). Associated with:
\begin{itemize}
    \item Kinship-based cooperation
    \item Personal exchange
    \item Group loyalty
\end{itemize}
\end{definition}

\subsection{Key Axioms}

\begin{axiom}[UNMAPPED_BE-1: Kinship-Morality Link]
\label{ax:be1}
Ancestral kinship tightness ($K$) predicts modern moral universalism ($U$):
\begin{equation}
U_i = \alpha + \beta \cdot K_{ancestors(i)} + \gamma \cdot X_i + \varepsilon_i
\end{equation}
where $\beta < 0$: tighter kinship $\rightarrow$ less universalism.
\textbf{EBF Connection:} Parameterizes CONTEXT-CULTURE for cooperation patterns.
\end{axiom}

\begin{axiom}[UNMAPPED_BE-2: Cognitive Uncertainty]
\label{ax:be2}
Agents form noisy percepts of decision-relevant variables:
\begin{equation}
\tilde{x} = x + \sigma \cdot \varepsilon, \quad \varepsilon \sim N(0, 1)
\end{equation}
This noise generates systematic biases through interaction with nonlinear utility.
\textbf{EBF Connection:} Extends CORE-AWARE with explicit noise parameter $\sigma$.
\end{axiom}

\begin{axiom}[UNMAPPED_BE-3: Correlation Neglect]
\label{ax:be3}
Agents underweight correlations in information sources:
\begin{equation}
\hat{\mu}_{posterior} = (1-\alpha) \cdot \mu_{prior} + \alpha \cdot \bar{s}
\end{equation}
where updating weight $\alpha$ is excessive relative to Bayesian benchmark when signals are correlated.
\textbf{EBF Connection:} Explains belief polarization in CORE-AWARE.
\end{axiom}

% ============================================================================
% COGNITIVE UNCERTAINTY
% ============================================================================
\section{Cognitive Uncertainty Framework}

\subsection{The Core Model}

Enke \& Graeber (2023) propose that many behavioral anomalies arise from a single source: \textbf{cognitive uncertainty} about one's own preferences and beliefs.

\begin{theorem}[Cognitive Uncertainty Effects]
When agents have noisy perception of values with uncertainty $\sigma$:
\begin{enumerate}
    \item \textbf{Risk aversion increases}: Uncertainty about outcomes amplifies caution
    \item \textbf{Present bias emerges}: Future values are noisier than present values
    \item \textbf{Probability weighting}: Small probabilities are overweighted due to representation noise
\end{enumerate}
\end{theorem}

\subsection{Formalization}

\begin{axiom}[UNMAPPED_BE-4: Noise-Rationality Tradeoff]
\label{ax:be4}
Behavioral ``biases'' decrease with:
\begin{enumerate}
    \item Higher stakes (increases attention, reduces $\sigma$)
    \item Lower complexity (reduces computational demands)
    \item More experience (improves calibration)
\end{enumerate}
\begin{equation}
\sigma = \sigma_0 \cdot f(complexity) \cdot g(stakes)^{-1} \cdot h(experience)^{-1}
\end{equation}
\textbf{EBF Connection:} Links CORE-AWARE to behavioral predictions through noise.
\end{axiom}

\subsection{Complexity and Time Preferences}

Enke, Graeber \& Oprea (2024) show that apparent hyperbolic discounting emerges from complexity:

\begin{tcolorbox}[colback=blue!5!white, colframe=blue!75!black]
\textbf{Key Finding:} When intertemporal choices are made computationally simple:
\begin{itemize}
    \item Present bias decreases by 40-60\%
    \item Time preferences become more consistent
    \item Hyperbolicity largely disappears
\end{itemize}
This suggests ``hyperbolic discounting'' may be a computational artifact, not a fundamental preference.
\end{tcolorbox}

\begin{axiom}[UNMAPPED_BE-5: Complexity-Time Link]
\label{ax:be5}
Apparent present bias ($\beta$) is a function of choice complexity ($C$):
\begin{equation}
\beta_{apparent} = \beta_{true} + \phi \cdot C
\end{equation}
where $\phi > 0$ captures the complexity-induced bias toward immediate options.
\textbf{EBF Connection:} Refines CORE-WHEN (V) time discounting with complexity modifier.
\end{axiom}

% ============================================================================
% MORAL VALUES AND BEHAVIOR
% ============================================================================
\section{Moral Values in Economic Contexts}

\subsection{Moral Values and Voting}

Enke (2020) demonstrates that moral values predict political behavior:

\begin{itemize}
    \item \textbf{Universalists}: Support redistribution, immigration, globalization
    \item \textbf{Communalists}: Support local community, tradition, in-group protection
\end{itemize}

The moral universalism index explains 30-40\% of the variation in voting patterns across US counties.

\subsection{Market Experience and Fairness}

Enke (2023) shows that market experience shifts fairness judgments:

\begin{axiom}[UNMAPPED_BE-6: Market Reference Points]
\label{ax:be6}
Market experience ($M$) shifts fairness reference points:
\begin{equation}
Fair_{price} = \alpha \cdot Cost + (1-\alpha) \cdot Market_{price}
\end{equation}
where $\alpha$ decreases with market experience.
\textbf{EBF Connection:} Refines CORE-WHAT (C) fairness dimension with context.
\end{axiom}

% ============================================================================
% GLOBAL PREFERENCES
% ============================================================================
\section{Global Evidence on Preferences}

\subsection{The GPS Dataset}

Falk, Becker, Dohmen, Enke, Huffman \& Sunde (2018) created the Global Preferences Survey:

\begin{itemize}
    \item 80,000+ respondents
    \item 76 countries
    \item 6 core preferences: risk, time, trust, altruism, reciprocity, patience
\end{itemize}

\subsection{Key Findings}

\begin{tcolorbox}[colback=green!5!white, colframe=green!75!black, title=GPS Key Results]
\begin{enumerate}
    \item \textbf{Large cross-country variation}: Preferences vary more between than within countries
    \item \textbf{Correlation structure}: Patience and risk tolerance are positively correlated
    \item \textbf{Development link}: Patience predicts economic development (r = 0.65)
    \item \textbf{Cultural persistence}: Preferences correlate with historical/ancestral variables
\end{enumerate}
\end{tcolorbox}

% ============================================================================
% EBF INTEGRATION
% ============================================================================
\section{EBF Integration}

\subsection{10C Mapping}

\begin{table}[h]
\centering
\begin{tabular}{lll}
\textbf{10C Element} & \textbf{Enke Contribution} & \textbf{Parameter} \\
\hline
CORE-WHO (AAA) & Moral circle boundaries & Universalism Index \\
CORE-WHAT (C) & Moral preferences dimension & $U_{moral}$ \\
CORE-HOW (B) & Culture-cooperation link & $\gamma_{kinship}$ \\
CORE-WHEN (V) & Complexity-time interaction & $\phi$ complexity factor \\
CORE-WHERE (BBB) & GPS global parameters & 76-country database \\
CORE-AWARE (AU) & Cognitive uncertainty & $\sigma$ noise parameter \\
CORE-READY (AV) & Stakes-attention link & $g(stakes)$ function \\
CORE-STAGE (AW) & Cultural transmission & Ancestral persistence \\
\end{tabular}
\caption{Enke contributions mapped to 10C framework}
\end{table}

\subsection{Key Parameter Values for BBB}

\begin{tcolorbox}[colback=green!5!white, colframe=green!75!black, title=Enke Parameters for BBB Repository]
\begin{itemize}
    \item $\sigma_{cognitive}$: 0.15-0.35 - Cognitive uncertainty parameter
    \item $\phi_{complexity}$: 0.4-0.6 - Complexity-time elasticity
    \item $\beta_{kinship}$: -0.25 - Kinship tightness effect on universalism
    \item $\alpha_{correlation}$: 0.3-0.5 - Degree of correlation neglect
    \item Patience-development $r$: 0.65 - Cross-country correlation
\end{itemize}
\end{tcolorbox}

% ============================================================================
% WORKED EXAMPLE
% ============================================================================
\section{Worked Example: Predicting Cooperation Patterns}

Consider predicting cooperation in a public goods game across two societies:
\begin{itemize}
    \item Society A: Tight kinship, low universalism (U = 0.3)
    \item Society B: Loose kinship, high universalism (U = 0.7)
\end{itemize}

\textbf{Prediction using Enke framework:}

\begin{align}
Coop_{in-group}^A &= baseline + \beta_{ingroup} \cdot (1-U_A) = 0.4 + 0.5 \cdot 0.7 = 0.75 \\
Coop_{out-group}^A &= baseline + \beta_{outgroup} \cdot U_A = 0.2 + 0.4 \cdot 0.3 = 0.32 \\
Coop_{in-group}^B &= 0.4 + 0.5 \cdot 0.3 = 0.55 \\
Coop_{out-group}^B &= 0.2 + 0.4 \cdot 0.7 = 0.48
\end{align}

\textbf{Result:} Society A shows high in-group/out-group differential (0.75 vs 0.32), while Society B shows more universal cooperation (0.55 vs 0.48).

% ============================================================================
% SUMMARY
% ============================================================================

% === AUTO-GENERATED ENTRIES (2026-02-22) ===

%%%%%%%%%%%%%%%%%%%%%%%%%%%%%%%%%%%%%%%%%%%%%%%%%%%%%%%%%%%%%%%%%%%%%%%%%%%%%%%
\subsubsection{2023: Market Experience is a Reference Point in Judgments of Fairn...}
%%%%%%%%%%%%%%%%%%%%%%%%%%%%%%%%%%%%%%%%%%%%%%%%%%%%%%%%%%%%%%%%%%%%%%%%%%%%%%%

\paragraph{Citation.}
Enke, Benjamin (2023). ``Market Experience is a Reference Point in Judgments of Fairness.'' \textit{Review of Financial Studies}.

\paragraph{10C Integration.}
\begin{itemize}[nosep]
  \item \textbf{Domain}: [To be specified]
  \item \textbf{use\_for}: LIT-ENK, LIT-ENKE
\end{itemize}


%%%%%%%%%%%%%%%%%%%%%%%%%%%%%%%%%%%%%%%%%%%%%%%%%%%%%%%%%%%%%%%%%%%%%%%%%%%%%%%
\subsubsection{2024: Associative Memory, Beliefs and Market Outcomes}
%%%%%%%%%%%%%%%%%%%%%%%%%%%%%%%%%%%%%%%%%%%%%%%%%%%%%%%%%%%%%%%%%%%%%%%%%%%%%%%

\paragraph{Citation.}
Enke, Benjamin and Schwerter, Frederik and Zimmermann, Florian (2024). ``Associative Memory, Beliefs and Market Outcomes.'' \textit{Journal of Political Economy}, 132.

\paragraph{10C Integration.}
\begin{itemize}[nosep]
  \item \textbf{Domain}: [To be specified]
  \item \textbf{use\_for}: LIT-ENK, LIT-ENKE
\end{itemize}


%%%%%%%%%%%%%%%%%%%%%%%%%%%%%%%%%%%%%%%%%%%%%%%%%%%%%%%%%%%%%%%%%%%%%%%%%%%%%%%
\subsubsection{2022: Moral Boundaries}
%%%%%%%%%%%%%%%%%%%%%%%%%%%%%%%%%%%%%%%%%%%%%%%%%%%%%%%%%%%%%%%%%%%%%%%%%%%%%%%

\paragraph{Citation.}
Enke, Benjamin (2022). ``Moral Boundaries.'' \textit{Working Paper}.

\paragraph{10C Integration.}
\begin{itemize}[nosep]
  \item \textbf{Domain}: [To be specified]
  \item \textbf{use\_for}: LIT-ENK, LIT-ENKE
\end{itemize}


%%%%%%%%%%%%%%%%%%%%%%%%%%%%%%%%%%%%%%%%%%%%%%%%%%%%%%%%%%%%%%%%%%%%%%%%%%%%%%%
\subsubsection{2022: Inflation and Cognitive Noise}
%%%%%%%%%%%%%%%%%%%%%%%%%%%%%%%%%%%%%%%%%%%%%%%%%%%%%%%%%%%%%%%%%%%%%%%%%%%%%%%

\paragraph{Citation.}
Enke, Benjamin and Kozlowski, Julian and Chahrour, Ryan (2022). ``Inflation and Cognitive Noise.'' \textit{Working Paper}.

\paragraph{10C Integration.}
\begin{itemize}[nosep]
  \item \textbf{Domain}: [To be specified]
  \item \textbf{use\_for}: LIT-ENK, LIT-ENKE
\end{itemize}


%%%%%%%%%%%%%%%%%%%%%%%%%%%%%%%%%%%%%%%%%%%%%%%%%%%%%%%%%%%%%%%%%%%%%%%%%%%%%%%
\subsubsection{2023: Confidence, Self-Selection, and Bias in the Aggregate}
%%%%%%%%%%%%%%%%%%%%%%%%%%%%%%%%%%%%%%%%%%%%%%%%%%%%%%%%%%%%%%%%%%%%%%%%%%%%%%%

\paragraph{Citation.}
Enke, Benjamin and Graeber, Thomas (2023). ``Confidence, Self-Selection, and Bias in the Aggregate.'' \textit{American Economic Review}, 113.

\paragraph{10C Integration.}
\begin{itemize}[nosep]
  \item \textbf{Domain}: [To be specified]
  \item \textbf{use\_for}: LIT-ENK, LIT-ENKE
\end{itemize}


%%%%%%%%%%%%%%%%%%%%%%%%%%%%%%%%%%%%%%%%%%%%%%%%%%%%%%%%%%%%%%%%%%%%%%%%%%%%%%%
\subsubsection{2024: The Cognitive Foundations of Mental Models}
%%%%%%%%%%%%%%%%%%%%%%%%%%%%%%%%%%%%%%%%%%%%%%%%%%%%%%%%%%%%%%%%%%%%%%%%%%%%%%%

\paragraph{Citation.}
Enke, Benjamin and Graeber, Thomas (2024). ``The Cognitive Foundations of Mental Models.'' \textit{Working Paper}.

\paragraph{10C Integration.}
\begin{itemize}[nosep]
  \item \textbf{Domain}: [To be specified]
  \item \textbf{use\_for}: LIT-ENK, LIT-ENKE
\end{itemize}


%%%%%%%%%%%%%%%%%%%%%%%%%%%%%%%%%%%%%%%%%%%%%%%%%%%%%%%%%%%%%%%%%%%%%%%%%%%%%%%
\subsubsection{2024: Quantifying Lottery Choice Complexity}
%%%%%%%%%%%%%%%%%%%%%%%%%%%%%%%%%%%%%%%%%%%%%%%%%%%%%%%%%%%%%%%%%%%%%%%%%%%%%%%

\paragraph{Citation.}
Enke, Benjamin and Shubatt, John (2024). ``Quantifying Lottery Choice Complexity.'' \textit{Working Paper}.

\paragraph{10C Integration.}
\begin{itemize}[nosep]
  \item \textbf{Domain}: [To be specified]
  \item \textbf{use\_for}: LIT-ENK, LIT-ENKE
\end{itemize}


%%%%%%%%%%%%%%%%%%%%%%%%%%%%%%%%%%%%%%%%%%%%%%%%%%%%%%%%%%%%%%%%%%%%%%%%%%%%%%%
\subsubsection{2025: Preferences Over Probabilities and Cognitive Imprecision}
%%%%%%%%%%%%%%%%%%%%%%%%%%%%%%%%%%%%%%%%%%%%%%%%%%%%%%%%%%%%%%%%%%%%%%%%%%%%%%%

\paragraph{Citation.}
Enke, Benjamin and Shubatt, John (2025). ``Preferences Over Probabilities and Cognitive Imprecision.'' \textit{Working Paper}.

\paragraph{10C Integration.}
\begin{itemize}[nosep]
  \item \textbf{Domain}: [To be specified]
  \item \textbf{use\_for}: LIT-ENK, LIT-ENKE
\end{itemize}

% === END AUTO-GENERATED ENTRIES ===
\section{Summary}

\begin{tcolorbox}[colback=gray!10!white, colframe=gray!75!black, title=Key Takeaways]
\begin{enumerate}
    \item \textbf{Moral Universalism:} Kinship structures persistently shape cooperation patterns
    \item \textbf{Cognitive Uncertainty:} Single noise parameter explains multiple biases
    \item \textbf{Correlation Neglect:} Systematic underweighting of information redundancy
    \item \textbf{Complexity-Time:} Apparent hyperbolic discounting is largely computational
    \item \textbf{Global Variation:} Preferences vary systematically across countries
    \item \textbf{EBF Integration:} Enke provides cultural Ψ parameters and cognitive noise foundations
\end{enumerate}
\end{tcolorbox}

% ============================================================================
% GLOSSARY
% ============================================================================
\section{Glossary}

See Appendix G (Glossary) for standard EBF terms. Enke-specific terms:

\begin{description}
    \item[Cognitive Uncertainty] Noise in perception/computation of decision-relevant quantities
    \item[Correlation Neglect] Tendency to treat correlated signals as independent
    \item[Kinship Tightness] Anthropological measure of family structure intensity
    \item[Moral Communalism] Ethical system with group-dependent obligations
    \item[Moral Universalism] Ethical system with universal, impartial obligations
\end{description}

% ============================================================================
% CRITICAL FOUNDATIONS
% ============================================================================
\section{Critical Foundations}

\begin{enumerate}
    \item Enke, B. (2019). Kinship, Cooperation, and the Evolution of Moral Systems. \textit{QJE}. [Moral Systems]
    \item Enke, B. \& Graeber, T. (2023). Cognitive Uncertainty. \textit{QJE}. [Core Theory]
    \item Enke, B. \& Zimmermann, F. (2019). Correlation Neglect in Belief Formation. \textit{REStud}. [Beliefs]
    \item Enke, B., Graeber, T. \& Oprea, R. (2024). Complexity and Time. \textit{JEEA}. [Time Preferences]
    \item Falk, A. et al. (2018). Global Evidence on Economic Preferences. \textit{QJE}. [GPS Data]
\end{enumerate}

% ============================================================================
% OPEN ISSUES
% ============================================================================
\section{Open Issues}

\begin{enumerate}
    \item \textbf{Noise vs. preferences:} How to distinguish cognitive noise from genuine preference heterogeneity?
    \item \textbf{Cultural change:} Can moral universalism be deliberately cultivated?
    \item \textbf{Institutional feedback:} How do institutions shape (not just reflect) moral values?
    \item \textbf{Complexity measurement:} How to objectively quantify decision complexity?
\end{enumerate}

% ============================================================================
% REFERENCES
% ============================================================================
\section{References}

\nocite{bcm_master}

For complete bibliography, see Master Bibliography (\texttt{bibliography/master\_references.bib}).

Key Enke references integrated in this appendix:
\begin{itemize}[nosep]
    \item enke2019moral, enke2020cognitive, enke2020whatyousee
    \item enke2023complexity, enke2020moral\_values, enke2022universalism
    \item enke2021boundaries, enke2017ancestral, enke2021market
    \item enke2019globalpref, enke2021ancestral\_env
\end{itemize}

\end{document}
